\begin{ouverture}[grandeurs]{10}{Longueurs et périmètres}{Partie 2 · Modéliser et mesurer}
{Unités · Contours · Disque et nombre pi · Figures composées}
\objectif{Choisir une unité et convertir une longueur.}
\objectif{Calculer le périmètre d’un carré, d’un rectangle et d’un disque.}
\objectif{Décomposer un contour et contrôler un calcul de longueur.}
\prerequis{Numération décimale ; périmètre comme longueur d’un contour.}
\end{ouverture}
\pagestyle{otMain}\essentiel\fontsize{11}{15}\selectfont
\phantomsection\label{cours-10-1}
\essSection{1}{Mesurer, estimer et convertir}
\begin{definition}Les unités de longueur sont construites à partir du mètre avec des préfixes : kilo- signifie mille, hecto- cent, déca- dix ; déci- un dixième, centi- un centième, milli- un millième. Ainsi 1 km $=1\,000$ m ; 1 hm $=100$ m ; 1 dam $=10$ m ; 1 dm $=0{,}1$ m ; 1 cm $=0{,}01$ m ; 1 mm $=0{,}001$ m.\end{definition}
\begin{savoirfaire}[Convertir en mètres]Deux unités voisines sont dans un rapport de 10 : 1 dm $=10$ cm et 1 cm $=\frac1{10}$ dm $=0{,}1$ dm. Pour convertir en mètres, on se demande combien de mètres vaut l’unité : $3{,}7$ km $=3{,}7\times1\,000$ m $=3\,700$ m ; 45 cm $=45\times0{,}01$ m $=0{,}45$ m.\end{savoirfaire}
\essSection{2}{Périmètres du carré et du rectangle}
\begin{definition}Le périmètre d’une figure plane est la longueur de son contour.\end{definition}
\begin{propriete}Carré de côté $c$ : $P=4\times c$. Rectangle de longueur $L$ et de largeur $\ell$ : $P=2\times(L+\ell)$. Pour un carré de côté $3{,}5$ cm, on remplace $c$ par $3{,}5$ : $P=4\times3{,}5=14$ cm.\end{propriete}
\begin{attention}Les longueurs doivent être exprimées dans la même unité : un rectangle de 2 m sur 80 cm a pour périmètre $2\times(2+0{,}8)=5{,}6$ m.\end{attention}
\clearpage
\phantomsection\label{cours-10-2}
\essSection{3}{Périmètre du disque et nombre pi}
\begin{propriete}Le périmètre d’un disque est proportionnel à son diamètre : on multiplie toujours le diamètre par le même nombre, noté $\pi$ (« pi »), qui vaut environ $3{,}14$. Si $d$ est le diamètre et $r$ le rayon : $P=\pi\times d=2\times\pi\times r$.\end{propriete}
\label{fig-cours-10-2-cours-cercle-diametre}\figureGeo{cours-cercle-diametre}

\begin{exemple}Disque de diamètre 10 cm : $P=\pi\times10$ cm, valeur exacte ; $P\approx3{,}14\times10=31{,}4$ cm, valeur approchée. Le nombre $\pi$ n’est pas décimal : $3{,}14$ en est seulement un arrondi.\end{exemple}
\essSection{4}{Contours de figures composées}
\begin{savoirfaire}[Suivre le contour]On parcourt le contour sans compter les segments intérieurs. Un rectangle de 6 cm sur 4 cm prolongé par un demi-disque de diamètre 4 cm a pour contour $6+4+6$ cm, plus un demi-périmètre de disque, $\pi\times4:2\approx6{,}28$ cm : environ $22{,}28$ cm.\end{savoirfaire}

\clearpage\rubrique{Découverte}
\begin{blocLivre}{ex-10-1}\begin{exercice}[1]{Convertir en mètres}
Convertir en mètres : $3{,}7$ km ; 45 cm ; 2 hm ; 8 mm ; 6 dam.
\end{exercice}\end{blocLivre}
\begin{blocLivre}{ex-10-2}\begin{exercice}[1]{Convertir depuis le mètre}
Convertir : $0{,}8$ m en cm ; $1\,250$ m en km ; $2{,}5$ m en mm ; 30 m en dam.
\end{exercice}\end{blocLivre}
\begin{blocLivre}{ex-10-3}\begin{exercice}[1]{Carré et rectangle}
Calculer le périmètre d’un carré de côté $3{,}5$ cm, puis celui d’un rectangle de $7{,}2$ m sur $4{,}5$ m.
\end{exercice}\end{blocLivre}
\begin{blocLivre}{ex-10-4}\begin{exercice}[1]{Périmètre d’un disque}
Calculer la valeur exacte, puis une valeur approchée avec $\pi\approx3{,}14$, du périmètre d’un disque de diamètre 10 cm, puis d’un disque de rayon 5 m.
\end{exercice}\end{blocLivre}
\begin{blocLivre}{ex-10-5}\begin{exercice}[1]{Reporter au compas}
Sans règle graduée, reporter au compas les trois côtés de la ligne brisée ABCD bout à bout sur la demi-droite, à partir de son origine O. Mesurer ensuite la longueur obtenue.
\label{fig-ex-10-5}\figureGeo{lon-5}
\end{exercice}\end{blocLivre}

\rubrique{Entraînement}
\begin{blocLivre}{ex-10-6}\begin{exercice}[2]{Une figure composée}
La figure est formée d’un rectangle de 6 cm sur 4 cm et d’un demi-disque de diamètre 4 cm. Calculer une valeur approchée de son périmètre.
\label{fig-ex-10-6}\figureGeo{lon-6}
\end{exercice}\end{blocLivre}
\begin{blocLivre}{ex-10-7}\begin{exercice}[2]{Le tour du stade}
Une piste est formée de deux lignes droites de 100 m et de deux demi-cercles de 60 m de diamètre. Calculer la longueur d’un tour, au dixième de mètre.
\label{fig-ex-10-7}\figureGeo{lon-7}
\end{exercice}\end{blocLivre}
\begin{blocLivre}{ex-10-8}\begin{exercice}[2]{Retrouver une dimension}
Un rectangle a pour périmètre 30 cm et pour longueur 9 cm. Quelle est sa largeur ? Un carré a pour périmètre 26 cm : quel est son côté ?
\end{exercice}\end{blocLivre}
\begin{blocLivre}{ex-10-9}\begin{exercice}[2]{Proportionnel au diamètre}
Recopier et compléter avec $\pi\approx3{,}14$. Que devient le périmètre quand le diamètre est multiplié par 3 ?\begin{center}\renewcommand{\arraystretch}{1.35}\begin{tabular}{c|ccccc}\toprule
Diamètre (cm) & 1 & 2 & 5 & 10 & 20\\\midrule
Périmètre (cm) & … & … & … & … & …\\
\bottomrule\end{tabular}\end{center}
\end{exercice}\end{blocLivre}
\begin{blocLivre}{ex-10-10}\begin{exercice}[2]{Le contour d’un L}
Les angles de la figure sont droits. Calculer les deux longueurs manquantes, puis le périmètre de la figure.
\label{fig-ex-10-10}\figureGeo{lon-10}
\end{exercice}\end{blocLivre}

\rubrique{Problèmes}
\begin{blocLivre}{ex-10-11}\begin{exercice}[3]{La roue du vélo}
La roue d’un vélo a un diamètre de 70 cm. Quelle distance parcourt le vélo quand la roue fait un tour ? Et 1 000 tours ? Répondre en mètres, puis en kilomètres pour 1 000 tours.
\end{exercice}\end{blocLivre}
\begin{blocLivre}{ex-10-12}\begin{exercice}[3]{La clôture}
Un jardin rectangulaire mesure $18{,}5$ m sur 12 m. On le clôture entièrement, sauf un portail de $1{,}5$ m. Le grillage coûte $2{,}40$ € le mètre. Quelle longueur de grillage faut-il ? Quel est son prix ?
\end{exercice}\end{blocLivre}
\begin{blocLivre}{ex-10-13}\begin{exercice}[3]{Le marathon}
Un marathon mesure $42{,}195$ km. Un tour de piste d’athlétisme mesure 400 m. Combien de tours complets faudrait-il faire, et quelle distance resterait-il ?
\end{exercice}\end{blocLivre}
\begin{blocLivre}{ex-10-14}\begin{exercice}[3]{Le rond-point}
Un rond-point est un cercle de 15 m de rayon. Une voiture en fait les trois quarts. Quelle distance parcourt-elle environ ?
\label{fig-ex-10-14}\figureGeo{lon-14}
\end{exercice}\end{blocLivre}
\begin{blocLivre}{ex-10-15}\begin{exercice}[3]{Vrai ou faux}
Dire si chaque phrase est vraie ou fausse, en justifiant.\begin{questions}
\item Si le côté d’un carré double, son périmètre double.
\item Deux rectangles de même périmètre sont superposables.
\end{questions}
\end{exercice}\end{blocLivre}

\clearpage
